%! AP Statistics — Unit 8, Lesson 8.4 — Student Packet Cover Sheet
\documentclass[10pt]{article}
\usepackage{apstats-article}
\usepackage{apstats-boxes}
\usepackage{tikz}
\usepackage{ltablex}
\keepXColumns

\begin{document}

\begin{tikzpicture}[remember picture, overlay]
  \fill[navy] (current page.north west) rectangle
    ([yshift=-2.2cm]current page.north east);
  \node[anchor=west, white, font=\large\bfseries] at
    ([xshift=1cm, yshift=-0.7cm]current page.north west)
    {\CourseName: \SchoolYear};
  \node[anchor=west, white, font=\normalsize] at
    ([xshift=1cm, yshift=-1.3cm]current page.north west)
    {Unit 8: Inference for Categorical Data: Chi-Square};
  \node[anchor=west, white, font=\normalsize\bfseries] at
    ([xshift=1cm, yshift=-1.9cm]current page.north west)
    {Lesson 8.4 \quad Expected Counts in Two-Way Tables};
\end{tikzpicture}

\vspace{2.4cm}
\namedateperiod

\begin{learningtargetbox}
\begin{itemize}
    \item I can calculate expected counts in a two-way table using the formula
          $\text{Expected} = \dfrac{(\text{row total}) \times (\text{column total})}{\text{table total}}$.
    \item I can explain what expected counts represent under $H_0$ for a chi-square test
          (what we would see if the two categorical variables were independent, or if the
          populations were homogeneous).
    \item I can check the large-counts condition for a two-way table by verifying that
          all \emph{expected} counts are at least 5.
\end{itemize}
\end{learningtargetbox}

\begin{tocbox}
\renewcommand{\arraystretch}{1.5}
\begin{tabularx}{\linewidth}{clXc}
    \textbf{\#} & \textbf{Component} & \textbf{Description} & \textbf{Score} \\
    \hline
    1 & Warm-Up        & Spiral review: chi-square GOF decision and conclusion       & \hfill/5 \\
    2 & Guided Notes   & Reading two-way tables; deriving and applying expected count formula & \hfill/10 \\
    3 & Group Activity & Computing expected counts; large-counts condition check     & \hfill/12 \\
    4 & Exit Ticket    & Independent calculation of expected counts in a 2$\times$3 table & \hfill/6 \\
    5 & Homework       & Practice problems on expected counts and the large-counts condition & \hfill/12 \\
    \hline
       & \multicolumn{2}{l}{\textbf{Total}} & \hfill/45 \\
\end{tabularx}
\end{tocbox}

\end{document}
